\documentclass[11pt,a4paper]{article}

\usepackage[spanish,es-noquoting]{babel}
\usepackage[utf8]{inputenc}
\usepackage[T1]{fontenc}
\usepackage{geometry}
\usepackage[hyphens,spaces,obeyspaces]{url}
\usepackage{hyperref}
\hypersetup{breaklinks=true} % Permite el salto de línea en los enlaces
\usepackage{xcolor}

\definecolor{raro-jwt-1}{RGB}{71, 61, 16}
\definecolor{raro-jwt-2}{RGB}{130, 69, 66}

\usepackage{adjustbox}
\usepackage{listings}
\lstdefinelanguage{JavaCustom}{
language=Java,
morekeywords={var}
}

\lstdefinelanguage{json}{
language=Java
}

\lstdefinelanguage{properties}{
  sensitive=false,
  morecomment=[l]{\#},
  morecomment=[l]{!},
  morestring=[b]"
}

% Definir el estilo y la sintaxis para YAML
\lstdefinelanguage{yaml}{
  keywords={true, false, null, y, n},
  keywordstyle=\color{blue}\bfseries,
  comment=[l]{\#},
  commentstyle=\color{gray}\ttfamily,
  stringstyle=\color{orange}\ttfamily,
  sensitive=false,
  morestring=[b]',
  morestring=[b]",
  alsoletter={-},
}

\lstset{
    inputencoding=utf8,
    extendedchars=true,
    literate=
        {á}{{\'a}}1
        {é}{{\'e}}1
        {í}{{\'i}}1
        {ó}{{\'o}}1
        {ú}{{\'u}}1
        {Á}{{\'A}}1
        {É}{{\'E}}1
        {Í}{{\'I}}1
        {Ó}{{\'O}}1
        {Ú}{{\'U}}1
        {ñ}{{\~n}}1
        {Ñ}{{\~N}}1,
basicstyle=\ttfamily\small,
keywordstyle=\color{blue}\bfseries,
stringstyle=\color{orange},
commentstyle=\color{gray}\itshape,
numberstyle=\tiny\color{gray},
identifierstyle=\color{black},
%backgroundcolor=\color{gray!10},
frame=none,
breaklines=true,
showstringspaces=false,
tabsize=1,
keepspaces=true,
columns=fullflexible,
numbers=none,
numbersep=8pt,
rulecolor=\color{black},
captionpos=b,
xleftmargin=1.5em,
framexleftmargin=1.5em
}

\usepackage{natbib}
\usepackage{graphicx}
\usepackage{endnotes}
\renewcommand{\notesname}{Notas finales}

\usepackage{tikz}
\usetikzlibrary{arrows.meta, positioning, backgrounds, fit, babel}
\pgfdeclarelayer{background}
\pgfsetlayers{background,main}

\usepackage{seqsplit}
\usepackage{float}

\geometry{margin=2.5cm}

\title{Gestión de \textit{Secretos} (contraseñas, rutas, etc.)\\AWS Secrets Manager}
\author{}
\date{\today}

\begin{document}

%\author{Oscar Dieste} % Your name

\date{\normalsize\today} % Today's date or a custom date

\maketitle

%\begin{abstract}
%\end{abstract}

{
\let\clearpage\relax
\tableofcontents
}

\section{Gestión manual de \textit{secretos}}

Por \textit{secreto}, debemos entender cualquier tipo de información que debe permanecer oculta, ya que su difusión implicaría algún riesgo, normalmente el acceso indebido. Los ejemplos típicos son las contraseñas que permiten acceder a ciertos servicios (como es el caso de las bases de datos), las API KEYs que permiten el acceso a APIs externas, tokens OAuth/M2M (machine-to-machine) utilizados por una aplicación para autenticarse ante ciertos microservicios, las claves privadas, etc.

Excluimos del concepto de secreto las claves de usuario. Desde luego, una clave de usuario debe mantenerse en secreto. Sin embargo, cuando una aplicación requiere una clave de usuario, el usuario puede proporcionarla introduciéndola directamente por teclado. Sin embargo, en el caso de los \textit{secretos}, la situación es distinta porque quien debe poseer ese secreto es algún tipo de aplicación, por ejemplo, el programa que necesita acceder a una base de datos. Esto exige que la clave de acceso esté disponible en algún lugar (no hay usuario que la pueda proporcionar), lo que propicia su acceso/interceptación y uso indebido.

Antes de contar con gestores centralizados (como es el caso de AWS \textit{Secrets Manager}, que introduciremos posteriormente), el almacenamiento de credenciales solía gestionarse mediante tres enfoques, cada uno con importantes vulnerabilidades:

\begin{itemize}
    \item \textbf{\textit{Hardcode} en el código fuente:} Los secretos (API KEYs, contraseñas, tokens, etc.) se incluyen directamente en el código, por ejemplo:

\begin{lstlisting}[language=Java,basicstyle=\ttfamily\footnotesize]
package com.ejemplo.db;

import java.sql.Connection;
import java.sql.DriverManager;
import java.sql.Exception;

public class DatabaseConnector {

    /*
    * Secretos
    */
    private static final String DB_URL = "jdbc:postgresql://db.ejemplo.com:5432/midatabase";
    private static final String DB_USER = "admin_user";
    private static final String DB_PASSWORD = "Password_1234";
    
    public static void main(String[] args) {
        DatabaseConnector app = DriverManager.getConnection(DB_URL, DB_USER, DB_PASSWORD);
        try (Connection conn = app.connect()) {
            ...
        } catch (SQLException e) {
            System.err.println("Error de conexión: " + e.getMessage());
        }
    }
}
\end{lstlisting}
    
    Este procedimiento es inseguro ya que cualquier persona con acceso al código fuente puede acceder al secreto. El secreto también podría recuperarse de una versión ejecutable, ya que los strings aparecen en el segmento de datos del código. 

    \item \textbf{Ficheros de configuración locales (\texttt{.env}, \texttt{.json}, \texttt{.yaml}):} Los secretos se guardan en ficheros (normalmente de texto) dentro del servidor o entorno local de cada desarrollador, como podría ser el siguiente \texttt{config.properties}:

\begin{lstlisting}[language=Properties,basicstyle=\ttfamily\footnotesize]
db.url=jdbc:postgresql://db.ejemplo.com:5432/midatabase
db.user=admin_user
db.password=Password_1234
\end{lstlisting}

Su uso sería, por ejemplo, el siguiente:

\begin{lstlisting}[language=Java, basicstyle=\ttfamily\footnotesize]
package com.ejemplo.db;

import java.io.InputStream;
import java.sql.Connection;
import java.sql.DriverManager;
import java.sql.Exception;
import java.util.Properties;

public class DatabaseConnector {

    public static void main(String[] args) {
        Properties properties = new Properties();
    
        // Cargar el archivo de propiedades
         desde el classpath
        try{
            InputStream input = Main.class.getClassLoader().getResourceAsStream("config.properties"));
            properties.load(input);
        } catch (Exception e) {
            System.err.println(e.getMessage());
            return;
        }
    
        // Extraer las credenciales y conectar
        String url = properties.getProperty("dbUrl");
        String user = properties.getProperty("dbUser");
        String password = properties.getProperty("dbPassword");
    
        DatabaseConnector app = DriverManager.getConnection(dbUrl, dbUser, dbPassword);
        try (Connection conn = app.connect()) {
            ...
        } catch (Exception e) {
            System.err.println(e.getMessage());
        }
    }
}
\end{lstlisting}

    Aunque no se suban al repositorio (por ejemplo, usando \texttt{.gitignore}), estos ficheros quedan dispersos en los equipos de desarrollo sin, probablemente, un control estricto de acceso. Además, los ficheros de configuración deben subirse a los servidores%
    \footnote{Subir los ficheros a los servidores introduce un problema adicional: 1) O bien dichos archivos se transfieren manualmente, por ejemplo, usando \texttt{sftp}, o bien 2) habría que utilizar un pipeline CI/CD local en alguna máquina de desarrollo.     
    Una alternativa intermedia sería confiar en pipelines CI/CD públicos con seguridad contrastada, tales como GitHub, Bitbucket o CodeDeploy, o un Jenkins institucional bien protegido. Sin embargo, en todos estos casos anteriores, volvemos a tener el problema de que los secretos estarían disponibles en los repositorios.}, %
    por lo que cualquiera que tenga acceso a los servidores podría obtener los secretos.

    \item \textbf{Variables de entorno} Definir las variables directamente en la configuración del servidor web o del contenedor (que es algo muy frecuente), por ejemplo:

\begin{lstlisting}[basicstyle=\ttfamily\footnotesize]

# Imagen base oficial de AWS para Java (Amazon Corretto 21)
FROM amazoncorretto:21-alpine

WORKDIR /app

ENV DB_URL="jdbc:postgresql://db.ejemplo.com:5432/midatabase" 
ENV DB_USER="admin_user"
ENV DB_PASSWORD="Password_1234"

COPY target/*.jar app.jar
EXPOSE 8080
ENTRYPOINT ["java", "-jar", "app.jar"]
\end{lstlisting}

\vspace{1em}

las cuales se podrían utilizar de la forma siguiente:

\vspace{1em}

\begin{lstlisting}[language=Java,basicstyle=\ttfamily\footnotesize]
package com.ejemplo.db;

import java.sql.Connection;
import java.sql.DriverManager;
import java.sql.Exception;

public class DatabaseConnector {

    /*
    * Secretos
    */
    String dbUrl = System.getenv("DB_URL");
    String dbUser = System.getenv("DB_USER");
    String dbPassword = System.getenv("DB_PASSWORD");
    
    public static void main(String[] args) {
        DatabaseConnector app = DriverManager.getConnection(dbUrl, dbUser, dbPassword);
        try (Connection conn = app.connect()) {
            ...
        } catch (SQLException e) {
            System.err.println("Error de conexión: " + e.getMessage());
        }
    }
}
\end{lstlisting}

Los problemas de las variables de entorno son muy parecidas a los problemas de los archivos de configuración.

\end{itemize}

\section{AWS Secrets Manager}

\textbf{AWS Secrets Manager} es un \textit{vault} (bóveda de seguridad) administrado por AWS y diseñado para almacenar, rotar, recuperar y auditar secretos a lo largo de su ciclo de vida, evitando la necesidad de mantenerlos estáticos en disco o en memoria. Existen otras soluciones en otras plataformas como GCP Secret Manager, Azure Key Vault, etc.

Lo más importante para nuestros propósitos es que el control de acceso a los secretos se gestiona mediante políticas de permisos nativas de AWS Identity and Access Management (IAM), limitando qué usuario o recurso (como una instancia EC2 o una función Lambda) puede descifrar cada secreto específico. Veremos esto inmediatamente a continuación.

\subsection{Creación de secretos mediante la consola web}

AWS Secrets Manager puede usarse directamente a través de la consola web de AWS. Sólo es necesario acceder al servicio y definir los secretos. La figura~\ref{fig:pantalla-secrets-manager} muestra el interfaz de usuario.

\begin{figure}
    \centering
    \includegraphics[width=\linewidth]{figs/pantalla-secrets-manager.png}
    \caption{Pantalla de Secrets manager que permite la creación de secretos}
    \label{fig:pantalla-secrets-manager}
\end{figure}

El interfaz de usuario sugiere que hay distintos tipos de secretos. Esto no es correcto. Un secreto es simplemente un objeto json formado por pares clave-valor. La pantalla de credenciales para \texttt{Amazon RDS} funciona únicamente como una plantilla que genera un objeto json que contiene campos predefinidos como \texttt{engine}, \texttt{host}, \texttt{username}, \texttt{password}, \texttt{dbname} y \texttt{port}. En la práctica, no es necesario respetar esta plantilla, \textbf{con la excepción de que se utilice la rotación de secretos}%
\endnote{La rotación de secretos es un proceso automatizado en AWS Secrets Manager que sustituye credenciales periódicamente. Se realiza en ocasiones automáticamente en AWS (como es el caso de las bases de datos Aurora) o, más frecuentemente, mediante función AWS Lambda orquestadora para actualizar simultáneamente la credencial en el recurso destino (como Aurora o RDS) y en Secrets Manager. No describiremos en este documento la rotación de secretos porque es un tema complejo que requiere su propia guía.
} . %
En ese caso, es conveniente respetar la plantilla porque AWS proporciona opciones por defecto para rotar los secretos que suponen que dicha plantilla ha sido aplicada. 


Utilizar la consola web es una muy buena opción en muchos casos. Un administrador podría definir los secretos, y estos nunca estarían disponibles en texto claro. Los secretos se podrían utilizar tal y como se indica en la sección~\ref{sec:uso}. Esto funciona bastante bien para bases de datos, API KEYs, y secretos que tienen permanencia en el tiempo. Para evitar el riesgo de filtrado, el administrador debería establecer la rotación de secretos.

\subsection{Generación de secretos mediante el CLI}

Utilizando el CLI, los secretos podrían aparecer explícitamente en los scripts, por ejemplo:

\begin{lstlisting}[language=bash,]
# Variables de configuracion (ejemplo generico)
DB_USER="admin"
DB_PASS="Un4 Contr4senn4 d|f|c|l;"
...
\end{lstlisting}

\vspace{1em}

El nombre del usuario privilegiado \textbf{depende del motor de base de datos}. En H2 (versión~3.0.0 de Event Hub) puede ser cualquiera; el laboratorio usa \texttt{admin}. En Aurora PostgreSQL (versión~4.0.0) el usuario \textit{master} es obligatorio y \textbf{no} puede llamarse \texttt{admin} (palabra reservada): allí el nombre elegido es \texttt{master}. El path del secreto en la~3.0.0 es \texttt{prod/h2/admin}; el sufijo \texttt{admin} nombra el rol en H2, no el usuario master de Aurora.

Pero esto supondría, de nuevo, que los secretos se almacenan en texto claro. En su lugar, hay que generar los secretos \textbf{automáticamente}. Tenemos tres opciones para ello:
\begin{enumerate}
    \item Utilizar \texttt{openssl} para generar una contraseña aleatoria y \texttt{jq} para construir la estructura JSON en memoria antes de enviarla a AWS Secrets Manager. Por ejemplo:
\begin{lstlisting}[language=bash,]
# Variables de configuracion (ejemplo con forma tipica de RDS/Aurora)
SECRET_NAME="prod/aurora/admin"

# En un ejemplo Aurora real (V4.0.0) el usuario master seria "master", no "admin".
DB_USER="admin"
DB_PASS=$(openssl rand -base64 32 | tr -dc 'a-zA-Z0-9' | head -c 16)
DB_HOST="eventhub-cluster-aurora.cluster-c1234567890.us-east-1.rds.amazonaws.com" # Endpoint del Cluster (Writer)
DB_NAME="eventhub"
DB_ENGINE="aurora-postgresql" # Usar "aurora-mysql" si es MySQL
DB_PORT=5432                  # Usar 3306 si es MySQL

DB_CLUSTER_ID="eventhub-cluster-aurora"

# Construir el JSON con la estructura estándar requerida para Amazon Aurora
SECRET_JSON=$(jq -n \
  --arg user "$DB_USER" \
  --arg pass "$DB_PASS" \
  --arg host "$DB_HOST" \
  --arg dbname "$DB_NAME" \
  --arg engine "$DB_ENGINE" \
  --arg cluster_id "$DB_CLUSTER_ID" \
  --argjson port "$DB_PORT" \
  '{
    username: $user,
    password: $pass,
    engine: $engine,
    host: $host,
    port: $port,
    dbname: $dbname,
    dbClusterIdentifier: $cluster_id
  }')
\end{lstlisting}

\vspace{1em}

\texttt{SECRET\_NAME} constituye el identificador único utilizado por AWS Secrets Manager para registrar, organizar y recuperar un secreto dentro de una cuenta y región determinadas. Se recomienda emplear la convención de ruta *nix (foo/bar) ya que es sencilla de entender y simplifica restringir permisos en AWS IAM mediante comodines (como \texttt{prod/aurora/*} o, en Event Hub~3.0.0, \texttt{prod/h2/*}).

Aunque no es relevante en este punto, conviene explicar qué significa \texttt{DB\_CLUSTER\_ID="eventhub-cluster-aurora"}. Las bases de datos AWS Aurora siempre se crean en un cluster, lo que permite redundancia, alto rendimiento y tolerancia a fallos de forma automática. También podría haber sido posible crear otra base de datos como una MySQL, para lo cual el cluster no habría sido necesario. En la versión~3.0.0 de Event Hub no hay clúster: la base es H2 embebida y el secreto (\texttt{prod/h2/admin}) solo guarda \texttt{username}, \texttt{password} y \texttt{dbname}.

Finalmente, \texttt{SECRET\_JSON} es el objeto json com pares clave-valor del que hemos hablado anteriormente. Aunque hemos creado \texttt{SECRET\_JSON}, todavía no lo hemos subido a Secrets Manager. Esto es, el JSON no ha salido de la máquina local. Pero antes de eso, las restantes opciones.

\item Usar \texttt{aws secretsmanager get-random-password} en lugar de \texttt{openssl}, de la forma siguiente:

\begin{lstlisting}[language=bash]
DB_PASS=$(aws secretsmanager get-random-password \
  --password-length 32 \
  --exclude-characters '/"@'\''\\' \
  --query RandomPassword --output text)
\end{lstlisting}

\item Permitir a Secrets Manager generar la credencial automáticamente, como se observa en este script:

\begin{lstlisting}[language=bash]
# Event Hub 3.0.0 (aws-scripts/secrets.sh): H2 embebida, sin host ni puerto.
SECRET_NAME="prod/h2/admin"

# El nombre del administrador depende de la base de datos.
# En H2 el nombre puede ser cualquiera (aqui, admin).
# En PostgreSQL/Aurora (V4.0.0) el usuario master es obligatorio (no puede llamarse admin).
DB_USER="admin"
DB_NAME="eventhub"

# La password la genera Secrets Manager (no sale del portatil).
# SecretStringTemplate: JSON sin password; GenerateStringKey: clave donde inserta AWS la password.
GENERATE_SPEC=$(jq -c -n \
  --arg user "$DB_USER" \
  --arg dbname "$DB_NAME" \
  '{
    SecretStringTemplate: ({
      username: $user,
      dbname: $dbname
    } | tostring),
    GenerateStringKey: "password",
    PasswordLength: 32,
    ExcludeCharacters: "/\"@'\''\\"
  }')
\end{lstlisting}

\vspace{1em}

El último comando es difícil de interpretar. Si ejecutamos el código anterior añadiendo \lstinline[language=bash]{echo $GENERATE_SPEC} al final, se obtendría algo equivalente a:
\begin{lstlisting}[language=json]
{"SecretStringTemplate":"{\"username\":\"admin\",\"dbname\":\"eventhub\"}","GenerateStringKey":"password","PasswordLength":32,"ExcludeCharacters":"/\"@'\\"}
\end{lstlisting}

\vspace{1em}

Lo que estamos haciendo es crear un objeto json con ciertas claves (como \textit{ExcludeCharacters}) que Secrets Manager conoce y puede interpretar a la hora de la generación de la password. Una de estas claves es \textit{SecretStringTemplate}, la cual debe contener lo que falta del secreto, esto es, todas las claves (\texttt{username}, \texttt{dbname}, \ldots) menos la password. Secrets Manager añadirá la password cuando cree el secreto (para crear el secreto ---hasta ahora sólo hemos creado un JSON---, ver la sección~\ref{sec:creacion-secreto}).

Si el secreto fuera para Aurora (versión~4.0.0), el JSON incluiría además \texttt{host}, \texttt{port}, \texttt{engine}, etc., y el \texttt{username} sería \texttt{master}, no \texttt{admin}. En esa versión, además, Aurora puede crear el secreto master al crear el clúster (\texttt{--manage-master-user-password}), como se explica en los apuntes de Aurora.

\end{enumerate}

Hay algunos \textbf{temas importantes que conviene destacar}:
\begin{itemize}
    \item Hemos utilizado como ejemplo los datos y credenciales de acceso a base de datos. Sin embargo, el mismo procedimiento serviría para cualquier tipo de secreto.

    \item Intuitivamente, la tercera de las opciones es muy atractiva porque la password nunca se almacena en el disco o la memoria local. Sin embargo, esta estrategia no es siempre la mejor (ver los dos puntos siguientes).

    \item Con ciertos elementos arquitectónicos, como las bases de datos, crear el secreto no es suficiente. El gestor de base de datos tiene sus propios usuarios. Entonces, después de crear el secreto, debemos hacer un \texttt{CREATE USER} (o un \texttt{ALTER USER} si estamos haciendo una actualización).

    \item Derivado de lo anterior, de poco nos sirve que AWS cree la clave automáticamente si posteriormente tenemos que leerla en el mismo script para realizar el \texttt{CREATE USER}. Para eso, casi mejor usar directamente la primera o segunda opciones%
    \endnote{Y todavía existe una cuarta opción, que describiremos en los apuntes de la base de datos Aurora: se trata de que la base de datos defina automáticamente la password durante su creación. En AWS esto tiene ciertas ventajas relacionadas con la rotación de claves.}.%
    
\end{itemize}

\subsection{Creación de los secretos usando el CLI}\label{sec:creacion-secreto}

Finalmente, nos queda crear, o subir si se quere, el secreto a Secrets Manager. En el primero de los casos anteriores, haríamos lo siguiente:

\begin{lstlisting}[language=bash]
# Crear el secreto en AWS Secrets Manager (opciones 1 y 2: el JSON ya incluye la password)
aws secretsmanager create-secret \
    --name "$SECRET_NAME" \
    --description "Credenciales de administracion ($DB_USER) para la base de datos '$DB_NAME'." \
    --secret-string "$SECRET_JSON"
\end{lstlisting}

\vspace{1em}

y en el tercero de los casos (el de Event Hub~3.0.0 / \texttt{secrets.sh}):

\begin{lstlisting}[language=bash]
# Crear el secreto: AWS genera la password y la inserta en el JSON
aws secretsmanager create-secret \
    --name "$SECRET_NAME" \
    --description "Credenciales de administracion ($DB_USER) para la base de datos '$DB_NAME'." \
    --generate-secret-string "$GENERATE_SPEC" \
    --query ARN \
    --output text
\end{lstlisting}

\vspace{1em}

La diferencia del tercer procedimiento es la presencia de \texttt{--generate-secret-string "\$GENERATE\_SPEC"}, que indica a Secrets Manager que tiene que generar la password automáticamente según esa especificación. Una ventaja es que la clave nunca llega a almacenarse en memoria local ni viaja por la conexión de AWS CLI: se genera dentro de AWS. El \texttt{--query ARN} evita imprimir el \texttt{SecretString} en la salida del script.

\subsection{Uso de los secretos}

Independientemente del modo en que los secretos han sido creados, los pasos para ser usados son los mismos en ambos casos. En primer lugar, hay que dar los permisos adecuados en IAM. 
Esto es un tema complejo para exponer aquí, pero está explicado en detalle en la \href{https://www.overleaf.com/read/tfyjkhsjhztf#20828d}{Guía de IAM}. Además, el ejemplo que se plantea en dicha guía es, precisamente, la asignación de permiso de lectura de secretos a una función Lambda y a una máquina EC2. 

Una vez asignados los permisos, el recurso que haya sido autorizado puede ejecutar la consulta sin necesidad de configurar claves. Por ejemplo, se podría consultar el secreto mediante el siguiente código Java (se muestra como un servicio de Spring Boot, pero no sería necesario hacerlo así; funcionaría exactamente igual, por ejemplo, dentro de una función \texttt{main()}):

\begin{lstlisting}[language=Java]
package com.example.demo.service;

import org.springframework.stereotype.Service;
import software.amazon.awssdk.regions.Region;
import software.amazon.awssdk.services.secretsmanager.SecretsManagerClient;
import software.amazon.awssdk.services.secretsmanager.model.GetSecretValueRequest;

@Service
public class SecretService {

    private final SecretsManagerClient secretsClient;

    public SecretService() {
        // Al no especificar credenciales explícitas, la SDK busca 
        // automáticamente el rol IAM asignado a la EC2 o Lambda.
        this.secretsClient = SecretsManagerClient.builder()
                .region(Region.US_EAST_1)
                .build();
    }

    // SecretsManagerException: La ejecución falla por falta de permisos IAM
    // o inexistencia del secreto
    // SdkClientException: La ejecución falla por problemas locales, normalmente
    // de conectividad a la red
    public String obtenerCredencialesDb() throws SecretsManagerException, SdkClientException {
        GetSecretValueRequest request = GetSecretValueRequest.builder()
                .secretId("prod/h2/admin")
                .build();

        return secretsClient.getSecretValue(request).secretString();
    }
}
\end{lstlisting}

\texttt{SecretsManagerClient} y \texttt{GetSecretValueReques} son clases del SDK de AWS que permiten acceder a los secretos. Para que este código funcione es necesario que se ejecute desde cualquier recurso de cómputo (compute service) de AWS, como podría ser una máquina EC2, a la cual se le hayan asignados los permisos adecuados. De lo contrario, el código fallará (como sugieren las excepciones que aparecen en el ejemplo anterior)

Otra alternativa oara acceder a los secretos cuando usamos Spring Boot es utilizar la librería \textit{starter} \texttt{spring-cloud-aws-starter-secrets-manager}, la cual nos permite consultar los secretos durante la inicialización e inyectarlos donde convenga. La aproximación es:

\begin{itemize}
    \item Añadir la dependencia Maven (o Gradle)

\begin{lstlisting}[language=XML]
<dependency>
    <groupId>io.awspring.cloud</groupId>
    <artifactId>spring-cloud-aws-starter-secrets-manager</artifactId>
</dependency>
\end{lstlisting}

La versión la fija el BOM \texttt{spring-cloud-aws-dependencies} (4.1.0 en Event Hub~3.0.0), compatible con Spring Boot~4.1. No hace falta fijar la región con \texttt{spring.cloud.aws.region.static}: el SDK usa la cadena por defecto (\texttt{AWS\_REGION}, perfil o metadatos de EC2).

\item Indicar la importación del secreto utilizando la opción \texttt{aws-secretsmanager:}:

\begin{lstlisting}[language=properties]
spring.config.import=classpath:cognito.properties,aws-secretsmanager:/prod/h2/admin
\end{lstlisting}

\noindent o en formato \texttt{.yaml}:

\begin{lstlisting}[language=yaml]
spring:
  config:
    import:
      - aws-secretsmanager:/prod/h2/admin
\end{lstlisting}

\item Inyectar los valores usando la anotación \texttt{@Value} donde corresponda, por ejemplo:

\begin{lstlisting}[language=Java]
public class DatabaseHealthService {

    @Value("${username}")
    private String dbUser;

    public void logConnectionUser() {
        // El valor proviene directamente del secreto /prod/h2/admin
        System.out.println("Usuario de BD configurado: " + dbUser);
    }
}
\end{lstlisting}
    
\end{itemize}

Es importante fijarse en \texttt{@Value("\$\{username\}")}. El secreto se importa en Spring Boot usando las mismas claves (username, password, engine, host, etc.) indicadas en el secreto.

Ahora bien, si el secreto utilizase las mismas claves estándar del framework (por ejemplo, \texttt{spring.datasource.username}), Spring Boot podría configurar la conexión a la base de datos de forma automática. Esto es interesante, pero exige cierto trabajo. Las claves de AWS son las que aparecen en el secreto que hemos usado de ejemplo, pero en Spring Boot las claves estándar son:

\begin{itemize}
\item \texttt{spring.datasource.url}: Cadena de conexión JDBC.
\item \texttt{spring.datasource.username}: Nombre del usuario.
\item \texttt{spring.datasource.password}: Contraseña.
\item \texttt{spring.datasource.driver-class-name}: Driver JDBC (\textit{opcional}, suele inferirse de la URL).
\end{itemize}

Si queremos que Spring Boot configure la conexión automáticamente, tenemos dos opciones: 1) usar las claves de Spring Boot en Secrets Manager (opción perfectamente viable) o 2) mapear las claves de AWS a claves de Spring Boot en el fichero \texttt{application.properties}, de la forma:

\begin{lstlisting}[language=properties]
# Event Hub 3.0.0: H2; username/password/dbname salen del secreto prod/h2/admin
spring.datasource.driverClassName=org.h2.Driver
spring.datasource.username=${username}
spring.datasource.password=${password}
spring.datasource.url=jdbc:h2:mem:${dbname};DB_CLOSE_DELAY=-1;DB_CLOSE_ON_EXIT=FALSE
\end{lstlisting}

En la versión~4.0.0 (Aurora) el mapeo es análogo, pero la URL es PostgreSQL (\texttt{jdbc:postgresql://\$\{host\}:\$\{port\}/\$\{dbname\}}) y el \texttt{username} del secreto master es \texttt{master}, no \texttt{admin}.

Una vez inicializada la base de datos, podría usarse normalmente utilizando JPA, o usando el bean \texttt{JdbcTemplate} en el caso de que se deseen realizar invocaciones directas en SQL.

\theendnotes

\end{document}
